\subsubsection{Anti-Landauer entropy and the index balance}

\begin{proposition}[Anti-Landauer entropy]\label{thm:anti-landauer-bits}\lean{Obsmetrologia.AntiLandauer.antiLandauerBits}\leanok\uses{thm:coordination-laplacian-dual}
Anti-Landauer entropy is the reversible harmonic side of the cover nerve:
$\check H^0=|C|$, the global-section count kept by settling.
\end{proposition}

\begin{proposition}[Euler characteristic of the cover nerve]\label{thm:nerve-euler-char}\lean{Obsmetrologia.AntiLandauer.eulerChar}\leanok\uses{thm:cover-nerve}
The Euler characteristic of the cover nerve is $\chi=|V|-|E|$ over
$\mathbb Z$.
\end{proposition}

\begin{theorem}[The index balance $\check H^0-\check H^1=\chi$]\label{thm:anti-landauer-index-balance}\lean{Obsmetrologia.AntiLandauer.index_balance}\leanok\uses{thm:anti-landauer-bits, thm:nerve-euler-char, thm:coordination-h1rank}
Anti-Landauer bits minus Landauer bits equals the Euler characteristic:
$\check H^0-\check H^1=\chi$. Settling preserves the index while distributing
it between harmonic content and obstruction content.
\end{theorem}

\begin{proposition}[Consensus and contradiction charges]\label{thm:consensus-contradiction-charges}\lean{Obsmetrologia.AntiLandauer.consensusCharge}\leanok\uses{thm:anti-landauer-bits, thm:coordination-h1rank}
The same index ledger can be read as two finite charges: consensus charge is
$\check H^0$, and contradiction charge is $\check H^1$.
\end{proposition}

\begin{theorem}[The signed charge ledger]\label{thm:net-charge-balance}\lean{Obsmetrologia.AntiLandauer.net_charge_balance}\leanok\uses{thm:anti-landauer-index-balance, thm:consensus-contradiction-charges}
Consensus charge minus contradiction charge is the conserved net charge
$\chi$.
\end{theorem}

\begin{theorem}[Matter and anti-matter are the signed contributions]\label{thm:matter-antimatter-balance}\lean{Obsmetrologia.AntiLandauer.matter_antimatter_balance}\leanok\uses{thm:net-charge-balance}
The Landauer side is the matter contribution $\check H^1$, the
anti-Landauer side is the anti-matter contribution $\check H^0$, and their
signed difference is the same net Euler charge.
\end{theorem}

\begin{theorem}[The zero-net vacuum]\label{thm:anti-landauer-vacuum}\lean{Obsmetrologia.AntiLandauer.vacuum_iff_charge_balance}\leanok\uses{thm:net-charge-balance}
The finite-index vacuum is exactly the state in which consensus and
contradiction charges cancel.
\end{theorem}

\begin{theorem}[Equal pair creation preserves the index]\label{thm:equal-pair-creation-preserves-net}\lean{Obsmetrologia.AntiLandauer.equal_pair_creation_preserves_net}\leanok\uses{thm:net-charge-balance}
Adding the same finite amount to the consensus side and to the contradiction
side preserves the net charge.
\end{theorem}

\begin{theorem}[Fixed-net cycle creation requires consensus creation]\label{thm:fixed-net-cycle-requires-consensus}\lean{Obsmetrologia.AntiLandauer.fixed_net_cycle_requires_consensus}\leanok\uses{thm:net-charge-balance}
At fixed net charge, increasing the contradiction side by $d$ forces the
consensus side to increase by the same $d$.
\end{theorem}

\begin{theorem}[Fixed-net consensus annihilation requires cycle annihilation]\label{thm:fixed-net-consensus-annihilation-requires-cycle-annihilation}\lean{Obsmetrologia.AntiLandauer.fixed_net_consensus_annihilation_requires_cycle_annihilation}\leanok\uses{thm:net-charge-balance}
At fixed net charge, decreasing the consensus side by $d$ forces the
contradiction side to decrease by the same $d$.
\end{theorem}

\begin{theorem}[Index-preserving steps have balanced cohomology deltas]\label{thm:preserves-euler-balanced-delta}\lean{Obsmetrologia.AntiLandauer.preservesEuler_iff_balanced_h0_h1_delta}\leanok\uses{thm:anti-landauer-index-balance}
A finite transition preserves $\chi$ exactly when the signed movement of
$\check H^0$ equals the signed movement of $\check H^1$.
\end{theorem}

\begin{theorem}[Closure index steps conserve net charge]\label{thm:closure-index-step-conserves-net}\lean{Obsmetrologia.AntiLandauer.ClosureIndexStep.preserves_matter_antimatter_balance}\leanok\uses{thm:preserves-euler-balanced-delta, thm:matter-antimatter-balance}
Any closure step projected to an index-preserving cover-nerve transition keeps
anti-matter minus matter fixed.
\end{theorem}

\begin{theorem}[Gauge Noether current for the index ledger]\label{thm:gauge-noether-current-index-ledger}\notready\uses{thm:ledger-noether, thm:closure-index-step-conserves-net, thm:percolation-gauge-measure, thm:mass-is-landauer}
The current finite theorem proves conservation of the index ledger. A genuine
Noether-current theorem still needs a carrier where the $D_4$ gauge action acts
on closure-index steps and where the conserved current is functorially read
from that symmetry action.
\end{theorem}

\begin{theorem}[Anti-descent and mirror temperature]\label{thm:anti-descent-mirror-temperature}\notready\uses{thm:fixed-net-cycle-requires-consensus, thm:positive-temperature-free-energy-flow, thm:stage-ledger, thm:ledger-noether}
The current finite theorem proves only the index-compensation law. A genuine
mirror temperature theorem must add a temporal or stochastic carrier showing
that the anti-descent builds $\check H^0$ while free-energy descent discharges
$\check H^1$.
\end{theorem}

\begin{theorem}[A reversible nerve's index is all harmonic]\label{thm:reversible-index-harmonic}\lean{Obsmetrologia.AntiLandauer.reversible_index_all_harmonic}\leanok\uses{thm:anti-landauer-index-balance}
When the Landauer side vanishes, the Euler characteristic is carried wholly
by the anti-Landauer harmonic side.
\end{theorem}

\begin{theorem}[The tree index balance]\label{thm:tree-index-balance}\lean{Obsmetrologia.AntiLandauer.tree_index_balance}\leanok\uses{thm:anti-landauer-index-balance, thm:tree-coordination-free}
For a finite tree nerve, $\check H^0=1$, $\check H^1=0$, and $\chi=1$.
This is the index balance on the connected acyclic case.
\end{theorem}

\begin{theorem}[The distinction-denomination functor]\label{thm:distinction-denomination-functor}\notready\uses{thm:anti-landauer-index-balance, thm:net-charge-balance, thm:settles-by-three, thm:config-burgers-additive}
\emph{(Open frontier --- stated, not formalized.)} A single denomination of lawful change
across the library's move carriers: a functor from a category whose morphisms are the admitted
moves --- Hoare traversals, Chu transforms, settling steps, defect slides, cover-nerve closure
steps --- to the signed charge ledger of \cref{thm:net-charge-balance}, assigning each move its
distinction count (collapsed distinctions on the $\check H^1$ side, kept sections on the
$\check H^0$ side), additive under composition, and conserved up to the Euler index of
\cref{thm:anti-landauer-index-balance}. The green pieces are the per-carrier accountings:
additivity of defect content (\cref{thm:config-burgers-additive}), stabilization of
knowledge-inflationary maps by the third iterate (\cref{thm:settles-by-three}), and the index
balance itself. Open: the common category of moves, the functor, and the proof that each
existing per-carrier accounting is its restriction. A countermodel --- a lawful move carrier
whose accounting cannot factor through one ledger --- would be equally decisive.
\end{theorem}
